\clearpage
\section{Dressed Schur returns and the exact positive prime norm}
\label{sec:dressed}
\subsection{The elementary representation and electric preparation}
Fix one real parameter $0<q<1$. This parameter will be the same for
all primes. The elementary $q$-Weyl representation associated with
SQED with one hypermultiplet has the following data
\cite[Section 3.1.1]{RG}. Let
\begin{equation}\label{eq:schur-space}
 \HH=\ell^2(\Z)\otimes L^2(S^1,\dd\theta/(2\pi)),
 \qquad \zeta=e^{i\theta},
\end{equation}
and write the magnetic charge as $m$. On the appropriate algebraic
line-state domain the actions are
\begin{align}
 (v\psi)_m(\zeta)&=q^{-m}\zeta\psi_m(\zeta),\notag\\
 (u_-\psi)_m(\zeta)&=(1+q^{-m-1}\zeta)\psi_{m+1}(q^{-1}\zeta),
 \notag\\
 (u_+\psi)_m(\zeta)&=\psi_{m-1}(q\zeta),\notag\\
 (\Omega_q)_m(\zeta)&=C_q\delta_{m0}\chi_q(\zeta),
 \qquad C_q=(q^2;q^2)_\infty.
 \label{eq:schur-actions}
\end{align}
The physical adjoints are inherited from the doubled representation;
we do not use $u_-=u_+^\dagger$. Positivity in the calculation is
the ordinary positive norm of \eqref{eq:schur-space}.
The elementary representation is a protected-sector control, not an
identification with the conformal four-flavour theory of
Section~\ref{sec:schur}.

For $0<a<1$ define, on the zero-charge sector,
\begin{equation}\label{eq:dressing}
 D_{q,a}(\zeta)=\frac{\sqrt{1-a^2}}{C_q}
 \frac{(-q\zeta;q^2)_\infty}{1+a\zeta}.
\end{equation}
The numerator is entire in $\zeta$ and the expression is analytic in
a neighbourhood of the closed unit disk. Thus its Taylor series
defines an electric-line completion, bounded on that sector.
It is not asserted to be bounded on all magnetic sectors. The prepared
state is explicitly
\begin{equation}\label{eq:eta}
 \eta_a=D_{q,a}(v)\Omega_q
 =\delta_{m0}\frac{\sqrt{1-a^2}}{1+a\zeta},\qquad
 \norm{\eta_a}=1.
\end{equation}
Multiplication by $-v$ on this sector gives
\begin{equation}\label{eq:eta-returns}
 \ip{\eta_a}{(-v)^k\eta_a}=a^{|k|},\qquad k\in\Z.
\end{equation}
Both assertions follow from the geometric coefficients
$\sqrt{1-a^2}(-a)^n$, $n\geq0$. Thus the descendant factors in
\eqref{eq:bare-overlap} have been removed by changing the source,
without changing the metric. The parameter $a$ is independent of $q$.
This infinite electric preparation is additional specified observable
data; it has not been derived as a unique UV boundary condition.

\subsection{One magnetic insertion and its graph domain}
\begin{lemma}\label{lem:one-magnetic}
If $a=qr$ with $0<r<1$, the vector $\eta_{qr}$ belongs to the graph
closure of the algebraic negative magnetic action on its electric-line
approximants. Its image is
\begin{equation}\label{eq:magnetic-state}
 (u_-\eta_{qr})_m(\zeta)
 =\delta_{m,-1}\sqrt{1-q^2r^2}\frac{1+\zeta}{1+r\zeta},
 \qquad
 \norm{u_-\eta_{qr}}^2=\frac{2(1-q^2r^2)}{1+r}.
\end{equation}
Every closed physical extension of that algebraic action contains this
graph limit and has the displayed value.
\end{lemma}
\begin{proof}
The formula for the state follows by substituting \eqref{eq:eta} in
\eqref{eq:schur-actions}. To justify that substitution as an operator
action, let $D_N$ be the Taylor polynomials of $D_{q,qr}$.
Its nearest possible pole has modulus $(qr)^{-1}>q^{-1}$, so choose
$R$ strictly between these two radii. The polynomials converge uniformly
on $|\zeta|\leq R$. Hence $D_N(v)\Omega_q\to\eta_{qr}$ in norm.
On the output sector,
\begin{equation}\label{eq:graph-cancellation}
 u_-D_N(v)\Omega_q
 =C_q\delta_{m,-1}D_N(q^{-1}\zeta)\chi_q(q\zeta),
\end{equation}
using $(1+\zeta)\chi_q(q^{-1}\zeta)=\chi_q(q\zeta)$.
This converges in norm to the displayed state. In particular no action
is being defined by crossing a pole and choosing a boundary value
afterwards. Finally,
\begin{equation}\label{eq:magnetic-coeff}
 \frac{1+\zeta}{1+r\zeta}
 =1+\sum_{n\geq1}(1-r)(-r)^{n-1}\zeta^n
\end{equation}
has squared norm $1+(1-r)^2/(1-r^2)=2/(1+r)$.
\end{proof}

\subsection{The full norm, including its contact}
For $d>0$, set
\begin{equation}\label{eq:beta}
 \beta_{q;r,d}^2=
 \frac{dr(1+r)}{(1-r)(1-q^2r^2)},\qquad
 \Phi_{q;r,d}[h]
 =\beta_{q;r,d}\,h(-q^{-1}v)u_-D_{q,qr}(v)\Omega_q
\end{equation}
for a finite Laurent polynomial $h$. On the output sector $m=-1$,
$v=q\zeta$, so $-q^{-1}v=-\zeta$. This normalization of the
electric argument is part of the source prescription.

\begin{theorem}[Schur realization of the positive local prime reference]
\label{thm:schur-prime}
With $z=-\zeta$, the preparation \eqref{eq:beta} satisfies
\begin{align}
 \ip{\Phi_{q;r,d}[h]}{\Phi_{q;r,d}[k]}
 &=\int_{S^1}\overline{h(z)}k(z)w_{r,d}(z)\frac{\dd\theta}{2\pi},
 \label{eq:full-schur-pairing}\\
 w_{r,d}(z)&=\frac{dr(1+r)}{1-r}\frac{|1-z|^2}{|1-rz|^2}
 =\kappa_{r,d}-d\sum_{n\geq1}r^n(z^n+z^{-n}),
 \label{eq:prime-weight}\\
 \kappa_{r,d}&=\frac{2dr}{1-r}.
 \label{eq:kappa}
\end{align}
The series converges absolutely and uniformly. In the Laurent basis,
\begin{equation}\label{eq:toeplitz}
 \ip{\Phi[z^j]}{\Phi[z^k]}=
 \begin{cases}
  \kappa_{r,d},&j=k,\\
  -dr^{|j-k|},&j\ne k.
 \end{cases}
\end{equation}
For each prime $p$, setting $r=p^{-1/2}$ and $d=\log p$ supplies every
single-prime off-diagonal coefficient in the Weil form. The same fixed
$q$ works for all these preparations.
\end{theorem}
\begin{proof}
Lemma~\ref{lem:one-magnetic} and \eqref{eq:beta} give the first expression
for $w_{r,d}$ directly from the physical norm. To expand it, use the
Poisson-kernel identity
\begin{equation}\label{eq:poisson}
 \frac{1-r^2}{|1-rz|^2}=1+\sum_{n\geq1}r^n(z^n+z^{-n})
 \qquad (|z|=1).
\end{equation}
Subtracting this identity from $(1+r)/(1-r)$ and multiplying by $d$
gives \eqref{eq:prime-weight}. Integration against
$\overline{z^j}z^k$ gives \eqref{eq:toeplitz}. Finally
$a_p=q/\sqrt p<q$, so the graph-domain proof applies to every prime
at the same $q$. All operators needed after the magnetic action are
bounded on its fixed output sector for Laurent polynomial inputs.
\end{proof}

The source normalization $\beta$ uses the assigned arithmetic boundary
data. It determines the entire multiplier; no independent contact is
subtracted. Since $w_{r,d}$ is bounded, \eqref{eq:full-schur-pairing}
extends by density from Laurent polynomials to $L^2(S^1)$ as a
bounded source map into the output sector.

\begin{corollary}[Sharp contact in this stationary scalar class]
\label{cor:contact}
A continuous scalar circle multiplier with nonzero Fourier coefficients
$-dr^{|n|}$ can be nonnegative only if its constant coefficient is at
least $\kappa_{r,d}$. The multiplier in \eqref{eq:prime-weight} attains
that value.
\end{corollary}
\begin{proof}
The multiplier with constant coefficient $C$ equals
$C-d\sum_{n\geq1}r^n(z^n+z^{-n})$. Its value at $z=1$ is
$C-2dr/(1-r)$, which must be nonnegative. Equation~\eqref{eq:prime-weight}
is nonnegative and vanishes at $z=1$.
\end{proof}
The same statement holds for the corresponding whole-line stationary
multiplier. It is not a claim of optimality over all finite-interval
or nonstationary realizations. This is the known positive local reference;
the result here supplies its Schur operator factor and checked domain.

\subsection{An explicit lift to continuous arithmetic inputs}
The circle coordinate is not, by itself, the real logarithmic input.
To give an actual source map, define
\begin{equation}\label{eq:bloch}
 (\mathcal B_dF)_j(\theta)=d^{-1/2}
 \widehat F((\theta+2\pi j)/d),\qquad
 0\leq\theta<2\pi,\quad j\in\Z.
\end{equation}
This is a unitary map from $L^2(\R)$ to
$L^2(S^1,\dd\theta/(2\pi);\ell^2(\Z))$. Indeed the intervals
$[2\pi j/d,2\pi(j+1)/d)$ partition the Fourier line, and
\eqref{eq:fourier} gives equality of norms. Conversely an arbitrary
square-integrable family specifies a Fourier transform interval by
interval, giving surjectivity.

Apply the extended source of Theorem~\ref{thm:schur-prime} to each
component of $\mathcal B_dF$ in an orthogonal copy of $\HH$. The
resulting positive norm is
\begin{align}
 B_{r,d}[F]
 &=\frac1{2\pi}\int_\R w_{r,d}(e^{id\tau})
                   |\widehat F(\tau)|^2\dd\tau\notag\\
 &=\kappa_{r,d}\norm F^2
   -2d\sum_{n\geq1}r^n\re\ip F{U_{nd}F}.
 \label{eq:lifted-norm}
\end{align}
The Fourier sign for $U_d$ does not change the expression, because
both orientations occur. For $F=E_Lf$, only $nd<L$ remains.

This construction adds a specified multiplicity space and preparation.
That multiplicity is not identified with the magnetic-charge label
already present in \eqref{eq:schur-space}. No particular four-dimensional
interface has yet been shown to implement \eqref{eq:bloch}. The result
is an exact positive Hilbert-space realization containing an elementary
Schur operator factor. A full field-theoretic interpretation of the
lift is part of the open matching problem.
